\section{Threats to validity}
\label{sec:threats}
\textbf{One benchmark, one split.} All results use MaliciousSkillBench's frozen \texttt{source\_disjoint}
split. The unseen benign test set has only \TestBenign{} skills from two sources (ATR and SkillFortifyBench),
and ATR contributes \TestAtrBenign{} of them, so benign FPR is largely an ATR number. A full
leave-one-source-out evaluation would show how much the conclusions depend on which sources are held out.

\textbf{``Mixed-only'' is effectively one source.} SkillTrojan contributes \MixedSkillTrojan{} skills and
AgentTrap \MixedAgentTrap{}; SkillTrustBench contributes the other \MixedSkillTrustBench{}. Our fix is therefore
best described as ``train on SkillTrustBench'', and its success may depend on how that source was built
(\S\ref{sec:mechanism}).

\textbf{Construction artifacts.} Inside SkillTrustBench most malicious skills are injected variants of benign
skills, and the held-out malicious sources are also injected or backdoored. A detector that learns ``find the
injected span'' may do well here and still miss malicious skills that were written from scratch.

\textbf{Thresholds.} Validation data comes from the training sources, so thresholds tuned on it do not
transfer to unseen sources (\S\ref{sec:fix}), as also reported for this benchmark by
SkillScreen~\cite{skillscreen2026}. Numbers labelled ``oracle'' pick the threshold on the test set and are
upper bounds.

\textbf{Prevalence.} The benchmark is 77\% malicious; real registries are mostly benign. Precision at
realistic prevalence would be much lower than the precision implied here, which is why we report FPR rather
than accuracy.

\textbf{Compute.} ModernBERT was fine-tuned once per training set (seed 0) on a laptop GPU, on the first 512
tokens only. Transformer results are a single run, not a distribution over seeds.

\section{Conclusion}
\label{sec:conclusion}
Malicious-skill detectors trained on merged benchmarks learn where a skill came from. Training only where both
labels share a source removes most of the resulting false positives, and a size-matched control rules out the
``less data'' explanation. \ConclusionMechanism{} For benchmark builders this suggests two concrete practices:
report benign FPR on unseen sources with the threshold fixed before looking at the test set, and document how
malicious and benign examples within each source were constructed, so that users can tell whether a detector
learned what makes a skill dangerous or how the benchmark was assembled.
